% ============================================================
% M3 S3 练习件·W2 臂D —— 课时17 2.8①压轴综合一（照抄课时01 母版体例）
% 生成：M3 成卷轮 S3 W2 臂D｜2026-09-14｜编译 cwd＝本目录（中文路径禁 \input 绝对路径）
%   xelatex main-true.tex ×2 → 含详解印本；xelatex main-false.tex ×2 → 纯题版（单源双本）
% 输入链（全只读）：题面库 课时17-2.8①压轴综合一.md（题面逐字装配）＋ -答案侧.md（值逐字回嵌）
%   ＋ 导学件17 main.tex 同键 ansblock（TeX 化详解底稿）；toolchain＝qp-m3.sty（钉后版
%   md5 7c3930362be8a0a2bdf21bbf8ac16573，本地挂载副本，破损即停）。
% 键账：件内 ansblock 键集 ≡ 题面库片练键集 ≡ manifest 练键序 01~16（16 键，零缺漏零多余）；
%   拓区隔离：2章-拓/测/滚 键不入本件（S4 拓展册收）；导学键 G1~G5 属导学件域，亦不入。
% 卷式例外案B：练习件不属卷式——答案块物理落点恒为题后 ansblock，禁卷末附卷制；
%   全线括线模（导言 \ansblockgrayfalse 一次置定，件内不切换；灰底盒不可跨栏断，练习件禁用）。
% 题序＝manifest 键序 01~16 零重排；印面题号 1~16 连号（\tihao 号＝\ansitem 号同键）；
%   三档切位 1~7／8~14／15~16（照库序切位，非难度重排）。
% 槽配实录：单选3（题4/5/8）＋填空4（题6/9/10/16）＋解答9（题1/2/3/7/11~15）；
%   照题面侧头标（批D §四 满足度表同口径）。
% 选项槽位：默认四连排（0.25\linewidth−0.25\lxhang），超宽降档梯 四连排→两连排
%   （0.5\linewidth−0.5\lxhang−0.25em）→单列 \lxopt；禁缩字号/负 kern/删标点凑宽。
%   本片实配：题4/题8 两连排、题5 单列（D 项「既不充分也不必要条件」超四连排档，导学件17
%   同款单列）（降档实测见值台账＋回执）。
% 解答书写区：九题 \liubai（带内 16~40mm）：题7 四判断 32mm、余默认 16mm；
%   纯题档吞 ansblock 不吞 \liubai。
% 填空印答：题6 \kongda{7}／题9 \kongda{16√3/5}／题10 \kongda{1}／题16 \kongda{1/3}
%   （详解版印答、纯题版退定宽空线，两档等形）。
% tieside：难度×知识点照同章导学件17 同题标签（逐题登记见值台账 难度词派生标注）。
% 印面清洗（导学件底稿→练习件印面）：剔 \ding 记号×10（题1/10/15/16 检验与互核句文句保留）；
%   跨题互引改印面题号：题2 互证句「变式 2、例 5」→「题5、题12」、题5「见例 2」→「见题2」、
%   题7(4)「与变式 2」→「与题5」、题12「此即例 2、变式 2」→「此即题2、题5」、
%   题14「正是例 4」→「正是题11」；题8 句尾元注记「（本席为点差法综合的铺垫位…）」剔；
%   题16 焦半径括注「到准线的距离，见课前预习」→「即到准线的距离」；题16 定位辨析句必带（承底稿）；
%   题面侧【注】行不入印面（题3 让位席注、题4 图嵌补全注）。
% ============================================================
\documentclass[fontset=none]{ctexart}
\usepackage{qp-m3}
% —— 印面逐字符号路由补钉（承导学件母版；− 走 TNR、▱ NSC 子块；禁改 sty 保 md5） ——
\xeCJKDeclareCharClass{Default}{"2212}%
\setCJKfamilyfont{nscblk}[Path=C:/提示词/工作区/字替对照-0909/variantF/fonts/,BoldFont=NSC-w600.ttf]{NSC-w500.ttf}%
\xeCJKDeclareSubCJKBlock{nscblk}{"25B1}%
\newcommand{\pxparallelogram}{{\CJKfamily{nscblk}▱}}%
% —— 断行微调（承导学件母版实测档，三零门承重；\emergencystretch 不另设） ——
\xeCJKsetup{CJKglue={\hskip 0.10em plus 0.08em minus 0.05em}}%
% —— 练习件全线括线模（S3 波次方案 §四.3：导言一次置定、件内不切换） ——
\ansblockgrayfalse
% —— 页脚件名（qp-headfoot 复用入口） ——
\renewcommand{\qpzhangming}{第二章\quad 平面解析几何}%
\renewcommand{\qpjianming}{练习件}%

\begin{document}
%装配轮S6三本跨本连续页码（G7方案B）setcounter 注入位
\setcounter{page}{230}

\keshi{第17课时\quad 2.8①压轴综合一}
\begin{multicols}{2}
\emergencystretch=1em

% ============================================================
% 基础巩固（题1~7＝17-01~17-07：解答6＋单选2＋填空1；题后紧跟答案制，全线括线 ansblock）
% ============================================================
\dangbadge{基}{础}{巩}{固}

\tihao{1}\zu{位置关系判定与公共点×解答} \tieside{简单(知识点一)}判断直线 \(y=-2x+4\) 与抛物线 \(y^{2}=4x\) 是否有公共点．如有，求出公共点的坐标，如公共点有两个，求出以这两个公共点为端点的线段长．
\liubai[16mm]{解答书写区}

\begin{ansblock}[2章-练-课时17-01]
% ans:2章-练-课时17-01
\ansitem{1}{有公共点；两公共点 (1,2)、(4,−4)；以之为端点的线段长 3√5}
\ansnote{详解}{以 \(y\) 为主元消元（直线已给 \(y\)、抛物线给 \(x=y^{2}/4\)）：\(y=-2\times(y^{2}/4)+4=-y^{2}/2+4\)，两边乘 2 整理得 \(y^{2}+2y-8=0\)（二次项系数 \(1\neq 0\)，故确为二次、不会退化）．\(\Delta=2^{2}+4\times 8=36>0\)，故有两个公共点．\(y=(-2\pm 6)/2\)，即 \(y=2\) 或 \(y=-4\)；代回 \(x=y^{2}/4\) 得 \(x=1\) 或 \(x=4\)．\par\noindent 故公共点为 \((1,2)\) 与 \((4,-4)\)，所求线段长 \(=\sqrt{(4-1)^{2}+(-4-2)^{2}}=\sqrt{9+36}=\sqrt{45}=3\sqrt{5}\)．（回验：\(2=-2\times 1+4\)；\(-4=-2\times 4+4\)；\(1=1\)、\(16=4\times 4\)．）}
\end{ansblock}

\tihao{2}\zu{相切\(\Leftrightarrow\)唯一公共点辨析×解答} \tieside{简单(知识点一)}举例说明，直线与抛物线只有一个公共点不是它们相切的充分条件．
\liubai[16mm]{解答书写区}

\begin{ansblock}[2章-练-课时17-02]
% ans:2章-练-课时17-02
\ansitem{2}{反例：直线 y=1 与抛物线 y²=4x 只有一个公共点 (1/4,1)，但二者不相切（该点处抛物线的切线是 y=2x+1/2）}
\ansnote{详解}{取抛物线 \(C\)：\(y^{2}=4x\) 与直线 \(y=1\)（平行于 \(C\) 的对称轴 \(x\) 轴）．把 \(y=1\) 代入得 \(1=4x\)，唯一解 \(x=1/4\)，故恰有一个公共点 \((1/4,1)\)．\par\noindent 再验它不相切：设过该点的直线 \(y=k(x-1/4)+1\)，记 \(t=x-1/4\)（则 \(4x=4t+1\)），代入 \(y^{2}=4x\) 得 \((1+kt)^{2}=1+4t\)，展开整理为 \(k^{2}t^{2}+(2k-4)t=0\)．\(k\neq 0\) 时方程有二根 \(t=0\) 与 \(t=(4-2k)/k^{2}\)，相切须二根重合，即 \((4-2k)/k^{2}=0\) \(\iff k=2\)，切线为 \(y=2(x-1/4)+1\)，即 \(y=2x+1/2\)；而 \(k=0\)（即直线 \(y=1\)）时联立式退化为一次方程 \(-4t=0\)，只有一个公共点却不是重根，按「相切＝联立有重根」的判定，它不是切线、直线与抛物线不相切．\par\noindent 故「直线与抛物线只有一个公共点」成立时相切未必成立，即它不是相切的充分条件．（本席事理与题5、题12 之退化特位三处互证．）}
\end{ansblock}

\tihao{3}\zu{弦长计算×解答} \tieside{简单(知识点三)}已知直线 \(x-2y+2=0\) 与椭圆 \(x^{2}+4y^{2}=4\) 相交于 \(A\)，\(B\) 两点，求线段 \(AB\) 的长．
\liubai[16mm]{解答书写区}

\begin{ansblock}[2章-练-课时17-03]
% ans:2章-练-课时17-03
\ansitem{3}{√5}
\ansnote{详解}{直线给出 \(x=2y-2\)，代入 \(x^{2}+4y^{2}=4\)：\((2y-2)^{2}+4y^{2}=4\to 4y^{2}-8y+4+4y^{2}=4\to 8y^{2}-8y=0\to 8y(y-1)=0\)，得 \(y=0\) 或 \(y=1\)，对应 \(x=-2\) 或 \(x=0\)．故 \(A(-2,0)\)、\(B(0,1)\)，\(|AB|=\sqrt{(0+2)^{2}+(1-0)^{2}}=\sqrt{5}\)．（此法直接解出两交点，无需弦长公式；若用弦长公式，\(x=my-2\) 型联立给 \(|y_{1}-y_{2}|=1\)、\(\sqrt{1+m^{2}}=\sqrt{5}\) 之倒数形式，结果同为 \(\sqrt{5}\)．）}
\end{ansblock}

\tihao{4}\zu{位置关系判定与公共点×单选} \tieside{简单(知识点一)}直线 \(y=x+1\) 与椭圆 \(x^{2}+y^{2}/2=1\) 的位置关系是\nobreak(\kongwei)

\bindopt {\fontsize{10.5pt}{\qpTJgLeadLiti}\selectfont \makebox[\dimexpr0.5\linewidth-0.5\lxhang-0.25em\relax][l]{A．相离}\makebox[\dimexpr0.5\linewidth-0.5\lxhang-0.25em\relax][l]{B．相切}\par}
\bindopt {\fontsize{10.5pt}{\qpTJgLeadLiti}\selectfont \makebox[\dimexpr0.5\linewidth-0.5\lxhang-0.25em\relax][l]{C．相交}\makebox[\dimexpr0.5\linewidth-0.5\lxhang-0.25em\relax][l]{D．无法确定}\par}

\begin{ansblock}[2章-练-课时17-04]
% ans:2章-练-课时17-04
\ansitem{4}{C}
\ansnote{详解}{联立 \(y=x+1\) 与 \(x^{2}+y^{2}/2=1\)，消去 \(y\)：\(x^{2}+(x+1)^{2}/2=1\)，乘 2 得 \(2x^{2}+x^{2}+2x+1=2\)，即 \(3x^{2}+2x-1=0\)．\(\Delta=2^{2}-4\times 3\times(-1)=4+12=16>0\)，且二次项系数 \(3\neq 0\)，故有两个公共点，直线与椭圆相交，选 C．（易错点：消元后须检查二次项系数是否为 0；本题 \(x^{2}\) 项由椭圆方程给出、恒在，无降次之虞．）}
\end{ansblock}

\tihao{5}\zu{相切\(\Leftrightarrow\)唯一公共点辨析×单选} \tieside{简单(知识点一)}「直线与抛物线相切」是「直线与抛物线只有一个公共点」的\nobreak(\kongwei)

\lxopt{A．充分不必要条件}

\lxopt{B．必要不充分条件}

\lxopt{C．充要条件}

\lxopt{D．既不充分也不必要条件}

\begin{ansblock}[2章-练-课时17-05]
% ans:2章-练-课时17-05
\ansitem{5}{A}
\ansnote{详解}{正向：相切即联立所得二次方程有重根（\(\Delta=0\)），重根只对应一个公共点，故「相切」\(\Rightarrow\)「只有一个公共点」，充分性成立．\par\noindent 反向：抛物线中消元可能降次——直线平行于对称轴时（如 \(y^{2}=4x\) 与 \(y=1\)，见题2），联立式退化为一次方程，恰有一个公共点却不相切，故「只有一个公共点」未必能推「相切」，必要性不成立．\par\noindent 于是「相切」是「只有一个公共点」的充分不必要条件，选 A．}
\end{ansblock}

\tihao{6}\zu{焦点弦×填空} \tieside{简单(知识点三)}已知直线 \(AB\) 过抛物线 \(y^{2}=4x\) 的焦点，交抛物线于 \(A(x_{1},y_{1})\)，\(B(x_{2},y_{2})\) 两点，若 \(x_{1}+x_{2}=5\)，则 \(|AB|\) 等于\kongda{7}．

\begin{ansblock}[2章-练-课时17-06]
% ans:2章-练-课时17-06
\ansitem{6}{7}
\ansnote{详解}{\(2p=4\)、\(p=2\)，焦点 \(F(1,0)\)、准线 \(x=-1\)．由抛物线定义，焦半径等于点到准线的距离：\(|AF|=x_{1}+1\)、\(|BF|=x_{2}+1\)．因 \(A\)、\(F\)、\(B\) 共线且 \(F\) 在 \(A\)、\(B\) 之间（焦点弦的两端点分居 \(x\) 轴两侧），\(|AB|=|AF|+|BF|=x_{1}+x_{2}+2=5+2=7\)．（即焦点弦公式 \(|AB|=x_{1}+x_{2}+p=5+2=7\)；按定义补出焦半径一步，防公式误记 \(p/2\)．）}
\end{ansblock}

\tihao{7}\zu{概念多判断×解答} \tieside{简单(知识点一)}判断正误．
(1) 过椭圆外一点只能作一条直线与椭圆相切．
(2) 直线 \(y=x\) 与椭圆 \(x^{2}/a^{2}+y^{2}/b^{2}=1\)（\(a>b>0\)）不一定相交．
(3) 过点 \((3,0)\) 的直线有且仅有一条与椭圆 \(x^{2}/9+y^{2}/16=1\) 相切．
(4) 直线与椭圆只有一个交点 \(\iff\) 直线与椭圆相切．
\liubai[32mm]{解答书写区}

\begin{ansblock}[2章-练-课时17-07]
% ans:2章-练-课时17-07
\ansitem{7}{(1)错误　(2)错误　(3)正确　(4)正确}
\ansnote{详解}{(1) 错．椭圆外一点可作两条切线：如过 \((4,0)\) 对椭圆 \(x^{2}/4+y^{2}=1\)，设 \(y=k(x-4)\) 代入得 \((1+4k^{2})x^{2}-32k^{2}x+64k^{2}-4=0\)，令 \(\Delta=16-192k^{2}=0\) 得 \(k=\pm\sqrt{3}/6\)，恰两条，故「只能一条」错．\par\noindent (2) 错．原点 \((0,0)\) 在椭圆内部（\(0+0<1\)），过内点的直线必与椭圆交于两点；直接联立 \(b^{2}x^{2}+a^{2}y^{2}=a^{2}b^{2}\) 与 \(y=x\) 得 \((a^{2}+b^{2})x^{2}=a^{2}b^{2}\)，即 \(x^{2}=a^{2}b^{2}/(a^{2}+b^{2})>0\)，确有两交点，故「不一定相交」错．\par\noindent (3) 对．\((3,0)\) 满足 \(9/9+0/16=1\)，是椭圆 \(x^{2}/9+y^{2}/16=1\) 的右顶点，在椭圆上；过椭圆上一点有且仅有一条切线（此处为 \(x=3\)），故对．\par\noindent (4) 对．椭圆是封闭曲线，直线与它联立消元后二次项系数恒不为 0（不存在像抛物线平行对称轴、双曲线平行渐近线那样的降次情形），故「恰一个公共点」\(\iff\)「方程有重根」\(\iff\)「相切」，与题5 中抛物线的反例恰成对照．\par\noindent 故四判断依次为「错、错、对、对」．}
\end{ansblock}

% ============================================================
% 综合提升（题8~14＝17-08~17-14：单选2＋填空3＋解答2，书写区 \liubai＋题后括线 ansblock）
% ============================================================
\dangbadge{综}{合}{提}{升}

\tihao{8}\zu{弦中点×单选} \tieside{简单(知识点三)}直线 \(y=x+1\) 被椭圆 \(x^{2}/4+y^{2}/2=1\) 所截得线段的中点的坐标是\nobreak(\kongwei)

\bindopt {\fontsize{10.5pt}{\qpTJgLeadLiti}\selectfont \makebox[\dimexpr0.5\linewidth-0.5\lxhang-0.25em\relax][l]{A．\((2/3,5/3)\)}\makebox[\dimexpr0.5\linewidth-0.5\lxhang-0.25em\relax][l]{B．\((4/3,7/3)\)}\par}
\bindopt {\fontsize{10.5pt}{\qpTJgLeadLiti}\selectfont \makebox[\dimexpr0.5\linewidth-0.5\lxhang-0.25em\relax][l]{C．\((-2/3,1/3)\)}\makebox[\dimexpr0.5\linewidth-0.5\lxhang-0.25em\relax][l]{D．\((-13/2,-17/2)\)}\par}

\begin{ansblock}[2章-练-课时17-08]
% ans:2章-练-课时17-08
\ansitem{8}{C}
\ansnote{详解}{联立消 \(y\)：\(x^{2}/4+(x+1)^{2}/2=1\)，乘 4 得 \(x^{2}+2(x+1)^{2}=4\)，即 \(3x^{2}+4x-2=0\)（\(\Delta=16-4\times 3\times(-2)=40>0\)，故确相交于两点）．设 \(A(x_{1},y_{1})\)、\(B(x_{2},y_{2})\)、中点 \(M(x_{0},y_{0})\)，韦达得 \(x_{1}+x_{2}=-4/3\)，故 \(x_{0}=(x_{1}+x_{2})/2=-2/3\)；中点在直线上，\(y_{0}=x_{0}+1=1/3\)．所以中点 \((-2/3,1/3)\)，选 C．}
\end{ansblock}

\tihao{9}\zu{弦长计算×填空} \tieside{简单(知识点三)}过双曲线 \(x^{2}/3-y^{2}/6=1\) 的右焦点作倾斜角为 \(30^{\circ}\) 的直线 \(l\)，直线 \(l\) 与双曲线交于不同的两点 \(A\)，\(B\)，则 \(AB\) 的长为\kongda{16√3/5}．

\begin{ansblock}[2章-练-课时17-09]
% ans:2章-练-课时17-09
\ansitem{9}{16√3/5}
\ansnote{详解}{\(a^{2}=3\)、\(b^{2}=6\)，\(c^{2}=a^{2}+b^{2}=9\)，右焦点 \(F_{2}(3,0)\)．倾斜角 \(30^{\circ}\) 得 \(k=\tan 30^{\circ}=\sqrt{3}/3\)，直线 \(l\)：\(y=(\sqrt{3}/3)(x-3)\)．代入 \(x^{2}/3-y^{2}/6=1\)：以 \(y^{2}=(x-3)^{2}/3\) 代，得 \(x^{2}/3-(x-3)^{2}/18=1\)，乘 18 整理为 \(6x^{2}-(x-3)^{2}=18\)，即 \(5x^{2}+6x-27=0\)（二次项系数 \(5\neq 0\)，故 \(l\) 不与渐近线平行；\(\Delta=36+540=576>0\)，确两交点）．韦达：\(x_{1}+x_{2}=-6/5\)、\(x_{1}x_{2}=-27/5\)，故 \((x_{1}-x_{2})^{2}=(x_{1}+x_{2})^{2}-4x_{1}x_{2}=36/25+108/5=576/25\)，\(|x_{1}-x_{2}|=24/5\)．弦长 \(|AB|=\sqrt{1+k^{2}}\,|x_{1}-x_{2}|=\sqrt{1+1/3}\times 24/5=(2/\sqrt{3})\times(24/5)=48/(5\sqrt{3})=16\sqrt{3}/5\)．}
\end{ansblock}

\tihao{10}\zu{弦长计数×填空} \tieside{简单(知识点二)}过双曲线 \(x^{2}/20-y^{2}/5=1\) 的右焦点的直线被双曲线所截得的弦长为 \(\sqrt{5}\)，这样的直线有\kongda{1}条．

\begin{ansblock}[2章-练-课时17-10]
% ans:2章-练-课时17-10
\ansitem{10}{1}
\ansnote{详解}{\(a^{2}=20\)、\(b^{2}=5\)，\(a=2\sqrt{5}\)、\(c^{2}=25\) 即右焦点 \(F(5,0)\)．过 \(F\) 的直线（除 \(x\) 轴外）可设为 \(x=my+5\)，代入 \(x^{2}-4y^{2}=20\) 得 \((my+5)^{2}-4y^{2}=20\)，整理为 \((m^{2}-4)y^{2}+10my+5=0\)．\par\noindent ① \(m^{2}=4\)（\(|m|=2\)，即直线斜率 \(|k|=1/2=b/a\)，平行渐近线）：二次项消失，化为一次方程 \(10my+5=0\)，只有一个公共点，不构成弦．\par\noindent ② \(m^{2}<4\)（斜率 \(|k|>1/2\)，较渐近线陡，两点同在右支）：\(y_{1}+y_{2}=-10m/(m^{2}-4)\)、\(y_{1}y_{2}=5/(m^{2}-4)\)，故 \((y_{1}-y_{2})^{2}=[100m^{2}-20(m^{2}-4)]/(m^{2}-4)^{2}=80(m^{2}+1)/(m^{2}-4)^{2}\)，\(|AB|=\sqrt{1+m^{2}}\,|y_{1}-y_{2}|=4\sqrt{5}(m^{2}+1)/(4-m^{2})\)．记 \(u=m^{2}\in[0,4)\)，\(f(u)=-1+5/(4-u)\) 随 \(u\) 单调增，即 \(|AB|\) 随 \(u\) 单调增，最小在 \(u=0\)，值 \(4\sqrt{5}\times(1/4)=\sqrt{5}\)，且仅此一处；\(u=0\) 即直线 \(x=5\)（垂直 \(x\) 轴的通径），其长 \(2b^{2}/a=2\times 5/(2\sqrt{5})=\sqrt{5}\)．\par\noindent ③ \(m^{2}>4\)（斜率 \(|k|<1/2\)，较渐近线平，两点分居两支）：\(|AB|=4\sqrt{5}(u+1)/(u-4)\)，而 \((u+1)/(u-4)=1+5/(u-4)\) 随 \(u\) 单调减，值域为 \((4\sqrt{5},+\infty)\)（\(u\to 4^{+}\) 时趋于 \(+\infty\)，\(u\to+\infty\) 时趋于 \(4\sqrt{5}\)＝实轴长但取不到，水平线 \(y=0\) 单列如下），恒大于 \(4\sqrt{5}>\sqrt{5}\)，取不到 \(\sqrt{5}\)．另单独核直线 \(y=0\)（\(x\) 轴，未被 \(x=my+5\) 覆盖）：它与双曲线交于两顶点，弦长 \(2a=4\sqrt{5}\neq\sqrt{5}\)．\par\noindent 综上，被截弦长等于 \(\sqrt{5}\) 的直线只有通径 \(x=5\) 这一条，故填 1．}
\end{ansblock}

\tihao{11}\zu{公共点个数分类求参×解答} \tieside{中档(知识点二)}已知双曲线 \(C\)：\(2x^{2}-y^{2}=2\) 与点 \(P(1,2)\)，讨论过点 \(P\) 的直线 \(l\) 的斜率的情况，使 \(l\) 与双曲线 \(C\) 分别有一个公共点、两个公共点、没有公共点．
\liubai[16mm]{解答书写区}

\begin{ansblock}[2章-练-课时17-11]
% ans:2章-练-课时17-11
\ansitem{11}{一个公共点：k∈\{√2,−√2,3/2\} 或 l 的斜率不存在；两个公共点：k∈(−∞,−√2)∪(−√2,√2)∪(√2,3/2)；没有公共点：k∈(3/2,+∞)}
\ansnote{详解}{① \(l\) 垂直 \(x\) 轴（斜率不存在）：直线 \(x=1\) 代入 \(2x^{2}-y^{2}=2\) 得 \(y=0\)，唯一公共点 \((1,0)\)（该处切线恰为 \(x=1\)），有一个公共点．\par\noindent ② \(l\) 不垂直 \(x\) 轴：设 \(y-2=k(x-1)\)，代入 \(2x^{2}-y^{2}=2\) 并整理得 \((2-k^{2})x^{2}+2(k^{2}-2k)x-k^{2}+4k-6=0\)．\par\noindent (a) \(2-k^{2}=0\)，即 \(k=\pm\sqrt{2}\)（\(l\) 平行渐近线 \(y=\pm\sqrt{2}x\)）：方程降为一次 \(2(k^{2}-2k)x-k^{2}+4k-6=0\)，一次项系数 \(2(2-2k)=4(1-k)\neq 0\)（\(k=\pm\sqrt{2}\) 时），故恰有一个公共点．\par\noindent (b) \(2-k^{2}\neq 0\)：\(\Delta=4(k^{2}-2k)^{2}+4(2-k^{2})(k^{2}-4k+6)=4(12-8k)=48-32k\)．\(\Delta=0\) \(\iff k=3/2\)，一个公共点（相切）；\(\Delta>0\) \(\iff k<3/2\)，两个公共点（本情形下须剔除 \(k=\pm\sqrt{2}\)，已归 (a)）；\(\Delta<0\) \(\iff k>3/2\)，没有公共点．\par\noindent 综上：有一个公共点 \(\iff k\in\{\sqrt{2},-\sqrt{2},3/2\}\) 或斜率不存在；有两个公共点 \(\iff k\in(-\infty,-\sqrt{2})\cup(-\sqrt{2},\sqrt{2})\cup(\sqrt{2},3/2)\)；没有公共点 \(\iff k\in(3/2,+\infty)\)．（注意 \(3/2\) 不在 \(k=\pm\sqrt{2}\) 之内（\(\sqrt{2}\approx 1.414<1.5\)），分段无冲突．）}
\end{ansblock}

\tihao{12}\zu{三态与特位×解答} \tieside{中档(知识点二)}设直线 \(l\)：\(y=kx+1\)，抛物线 \(C\)：\(y^{2}=4x\)，当 \(k\) 为何值时，\(l\) 与 \(C\) 相切？相交？相离？
\liubai[16mm]{解答书写区}

\begin{ansblock}[2章-练-课时17-12]
% ans:2章-练-课时17-12
\ansitem{12}{k=1 时相切；k<1 时相交；k>1 时相离}
\ansnote{详解}{联立 \(y=kx+1\) 与 \(y^{2}=4x\) 消 \(y\)：\((kx+1)^{2}=4x\)，即 \(k^{2}x^{2}+(2k-4)x+1=0\)．\par\noindent ① \(k\neq 0\)：二次项系数 \(k^{2}>0\)，为一元二次方程，\(\Delta=(2k-4)^{2}-4k^{2}=4k^{2}-16k+16-4k^{2}=16(1-k)\)．\(\Delta=0\) \(\iff k=1\)，相切；\(\Delta>0\) \(\iff k<1\) 且 \(k\neq 0\)，相交（两交点）；\(\Delta<0\) \(\iff k>1\)，相离．\par\noindent ② \(k=0\)：直线为 \(y=1\)，平行于对称轴 \(x\) 轴，代入得 \(1=4x\) 即 \(x=1/4\)，恰一个公共点 \((1/4,1)\)，属相交（不切）——此即题2、题5 的退化特位，并入「相交」档．\par\noindent 综上：\(k=1\) 相切；\(k<1\)（含 \(k=0\)）相交；\(k>1\) 相离．}
\end{ansblock}

\tihao{13}\zu{公共点个数分类求参×解答} \tieside{中档(知识点二)}已知直线 \(l\)：\(y=kx+2\) 和椭圆 \(C\)：\(x^{2}/3+y^{2}/2=1\)．分别求直线 \(l\) 与椭圆 \(C\) 有两个公共点、只有一个公共点和没有公共点时 \(k\) 的取值范围．
\liubai[16mm]{解答书写区}

\begin{ansblock}[2章-练-课时17-13]
% ans:2章-练-课时17-13
\ansitem{13}{两个公共点：k<−√6/3 或 k>√6/3；只有一个公共点：k=±√6/3；没有公共点：−√6/3<k<√6/3}
\ansnote{详解}{把 \(y=kx+2\) 代入 \(x^{2}/3+y^{2}/2=1\)，乘 6 得 \(2x^{2}+3(kx+2)^{2}=6\)，即 \((2+3k^{2})x^{2}+12kx+6=0\)．二次项系数 \(2+3k^{2}\geqslant 2>0\) 恒成立（椭圆无降次情形），故公共点个数完全由 \(\Delta\) 定：\(\Delta=(12k)^{2}-4(2+3k^{2})\times 6=144k^{2}-48-72k^{2}=72k^{2}-48\)．\par\noindent \(\Delta>0\) \(\iff k^{2}>48/72=2/3\) \(\iff|k|>\sqrt{6}/3\)，两个公共点；\(\Delta=0\) \(\iff k=\pm\sqrt{6}/3\)，只有一个公共点（相切）；\(\Delta<0\) \(\iff k^{2}<2/3\) \(\iff|k|<\sqrt{6}/3\)，没有公共点．（边界核对：直线恒过 \((0,2)\)，而 \((0,2)\) 在椭圆外（\(0+4/2=2>1\)），故存在相离与相切位置、与读数自洽．）}
\end{ansblock}

\tihao{14}\zu{三态与特位×解答} \tieside{中档(知识点二)}过原点的直线 \(l\) 与双曲线 \(x^{2}/4-y^{2}/3=1\) 相交于两点，求 \(l\) 的斜率的取值范围．
\liubai[16mm]{解答书写区}

\begin{ansblock}[2章-练-课时17-14]
% ans:2章-练-课时17-14
\ansitem{14}{|k|<√3/2（即 k∈(−√3/2, √3/2)）}
\ansnote{详解}{设 \(l\)：\(y=kx\)（过原点的直线中 \(x=0\) 即 \(y\) 轴与双曲线无公共点，故可只议斜率存在者）．代入 \(x^{2}/4-y^{2}/3=1\)，乘 12 得 \(3x^{2}-4k^{2}x^{2}=12\)，即 \((3-4k^{2})x^{2}=12\)．\par\noindent ① \(3-4k^{2}>0\)（\(|k|<\sqrt{3}/2\)）：\(x^{2}=12/(3-4k^{2})>0\)，解得 \(x=\pm\sqrt{12/(3-4k^{2})}\)，对应两个不同交点（关于原点对称），符合「相交于两点」．\par\noindent ② \(3-4k^{2}\leqslant 0\)（\(|k|\geqslant\sqrt{3}/2\)）：当 \(=0\)（\(l\) 与渐近线 \(y=\pm(\sqrt{3}/2)x\) 平行，且过原点即为渐近线本身）时方程为 \(0=12\) 无解；当 \(<0\) 时 \(x^{2}<0\) 无实数解；两种都无公共点．\par\noindent 故所求范围为 \(|k|<\sqrt{3}/2\)．（几何直观：过原点（双曲线中心）的直线只要比两条渐近线陡即与双曲线交于两支上对称两点；渐近线本身与曲线零公共点，正是题11 中 \(k=\pm\sqrt{2}\) 型降次的对照位．）}
\end{ansblock}

% ============================================================
% 思维探索（题15~16＝17-15~17-16：解答1＋填空1 难端补位对，居末档照库序）
% ============================================================
\dangbadge{思}{维}{探}{索}

\tihao{15}\zu{弦中点×解答} \tieside{难(知识点三)}已知椭圆 \(x^{2}/6+y^{2}/3=1\)，直线 \(x-y+1=0\) 与该椭圆交于 \(A\)，\(B\) 两点，求线段 \(AB\) 中点的坐标和线段 \(AB\) 的长．
\liubai[16mm]{解答书写区}

\begin{ansblock}[2章-练-课时17-15]
% ans:2章-练-课时17-15
\ansitem{15}{中点 (−2/3, 1/3)；|AB|=8√2/3}
\ansnote{详解}{直线写成 \(y=x+1\)，代入 \(x^{2}/6+y^{2}/3=1\)，乘 6 得 \(x^{2}+2(x+1)^{2}=6\)，即 \(3x^{2}+4x-4=0\)（\(\Delta=16+48=64>0\)，确两交点）．设 \(A(x_{1},y_{1})\)、\(B(x_{2},y_{2})\)，则 \(x_{1}+x_{2}=-4/3\)、\(x_{1}x_{2}=-4/3\)．\par\noindent 中点：\(x_{0}=(x_{1}+x_{2})/2=-2/3\)，\(y_{0}=x_{0}+1=1/3\)，即中点 \((-2/3,1/3)\)．\par\noindent 弦长：\((x_{1}-x_{2})^{2}=(x_{1}+x_{2})^{2}-4x_{1}x_{2}=16/9+16/3=64/9\)，\(|x_{1}-x_{2}|=8/3\)；直线斜率 \(k=1\)，故 \(|AB|=\sqrt{1+k^{2}}\,|x_{1}-x_{2}|=\sqrt{2}\times 8/3=8\sqrt{2}/3\)．\par\noindent（回代验：方程 \(3x^{2}+4x-4=0\) 之根 \(x=2/3\)、\(x=-2\)，对应点 \((2/3,5/3)\)、\((-2,-1)\)：中点 \(((2/3-2)/2,(5/3-1)/2)=(-2/3,1/3)\)，\(|AB|=\sqrt{(2/3+2)^{2}+(5/3+1)^{2}}=(8/3)\sqrt{2}\)．）}
\end{ansblock}

\tihao{16}\zu{焦点弦×填空} \tieside{难(知识点三)}过抛物线 \(x^{2}=2py\)（\(p>0\)）的焦点 \(F\) 作倾角为 \(30^{\circ}\) 的直线，与抛物线分别交于 \(A\)、\(B\) 两点（\(A\) 在 \(y\) 轴左侧），则 \(|AF|/|FB|=\)\kongda{1/3}．

\begin{ansblock}[2章-练-课时17-16]
% ans:2章-练-课时17-16
\ansitem{16}{1/3}
\ansnote{详解}{焦点 \(F(0,p/2)\)、准线 \(y=-p/2\)．倾角 \(30^{\circ}\) 即直线斜率 \(\tan 30^{\circ}=\sqrt{3}/3\)，直线 \(l\)：\(y=(\sqrt{3}/3)x+p/2\)，等价地 \(x=\sqrt{3}(y-p/2)\)．代入 \(x^{2}=2py\) 得 \(3(y-p/2)^{2}=2py\)，展开：\(3y^{2}-3py+3p^{2}/4-2py=0\)，即 \(3y^{2}-5py+3p^{2}/4=0\)，除以 3：\(y^{2}-(5p/3)y+p^{2}/4=0\)（\(\Delta=25p^{2}/9-p^{2}=16p^{2}/9\)，两根 \(y=(5p/3\pm 4p/3)/2\)，即 \(p/6\) 与 \(3p/2\)）．\par\noindent 为免 \(y\) 轴左右两侧混判，改以 \(x\) 为主元：把 \(y=x/\sqrt{3}+p/2\) 代 \(x^{2}=2py\) 得 \(x^{2}=2p(x/\sqrt{3}+p/2)=(2p/\sqrt{3})x+p^{2}\)，即 \(x^{2}-(2p/\sqrt{3})x-p^{2}=0\)．解之：\(x=(p/\sqrt{3})\pm(1/2)\sqrt{16p^{2}/3}=(p/\sqrt{3})\pm(2p/\sqrt{3})\)，故 \(x=3p/\sqrt{3}=\sqrt{3}p\) 或 \(x=-p/\sqrt{3}\)．\par\noindent 定位辨析：题设「\(A\) 在 \(y\) 轴左侧」即 \(x_{A}<0\)，故只能取负根 \(x_{A}=-p/\sqrt{3}\)（此时 \(y_{A}=x_{A}^{2}/(2p)=(p^{2}/3)/(2p)=p/6\)），而 \(x_{B}=\sqrt{3}p\)（\(y_{B}=3p^{2}/(2p)=3p/2\)）——此两步不可互换，否则比值倒置为 3．\par\noindent 焦半径（开口向上：\(|MF|=y_{M}+p/2\)，即到准线的距离）：\(|AF|=y_{A}+p/2=p/6+p/2=2p/3\)；\(|FB|=y_{B}+p/2=3p/2+p/2=2p\)．\par\noindent 故 \(|AF|/|FB|=(2p/3)/(2p)=1/3\)．（互核一（代回曲线）：\(A(-p/\sqrt{3},p/6)\) 使 \(x^{2}=p^{2}/3=2p\times p/6\)；\(B(\sqrt{3}p,3p/2)\) 使 \(x^{2}=3p^{2}=2p\times(3p/2)\)．互核二（焦点弦坐标积）：本题抛物线开口向上，焦点弦之对偶形式为 \(x_{1}x_{2}=-p^{2}\)，本处 \((-p/\sqrt{3})\times\sqrt{3}p=-p^{2}\)．互核三（弦长）：\(|AF|+|FB|=2p/3+2p=8p/3\)，而 \(|AB|=y_{A}+y_{B}+p=p/6+3p/2+p=8p/3\)，与两点距 \(\sqrt{(4p/\sqrt{3})^{2}+(4p/3)^{2}}=\sqrt{64p^{2}/9}=8p/3\) 合．）}
\end{ansblock}

% ---- 尾块（\tailfill 置 \end{multicols} 前；无参＝内容尾随框，每件恰 1 处） ----
\tailfill

\end{multicols}
\end{document}
